%2026v0
\documentclass{webofc}
\usepackage[varg]{txfonts}
\newcommand{\comment}[1]{}
\usepackage{graphicx}
\usepackage{listings}

\usepackage{booktabs}
\usepackage{multirow}
\usepackage{xcolor}
\usepackage{pifont}

% Custom macros matching the rst roles/replacements
\newcommand{\cmark}{\ding{51}} % Yes / Y checkmark or character
\newcommand{\xmark}{\ding{55}} % No / N mark or character
\newcommand{\btext}[1]{\textbf{#1}}
\newcommand{\rtext}[1]{\textbf{\color{red}#1}}


\lstset{
  basicstyle=\ttfamily\small,
  captionpos=b,                 % Places caption at the bottom (like figures)
  frame=single,                 % Adds a simple border around the listing
  numbers=left,                 % Line numbers on the left
  %numberstyle=\tiny\color{gray},
  mathescape=true               % Allows math mode like $\gets$ inside code
}


\usepackage{array}
\comment{

Sections from presentation : opticks\_20260527\_bangkok\_chep2026.txt

* Optical Photon Simulation Context

  * (JUNO) Optical Photon Simulation Problem...
  * Optical Photon Simulation ≈ Ray Traced Image Rendering
  * NVIDIA RTX PRO 6000 Blackwell
  * Four NVIDIA RTX Gen : Ray Tracing : ~2x every ~2 yrs
  * Opticks Optical simulation directly scales with ray tracing speed
  * NVIDIA OptiX : Ray Tracing Engine -- Accessible GPU Ray Tracing
  * Geant4 + Opticks + NVIDIA OptiX : Hybrid Workflow
  * Amdahl 'Law' : Overall Speedup S(p,n) for Parallel Fraction p and Parallized Speedup n

* Simulation of High Energy Events, eg: Multi-Muon

  * Photons from muon crossing JUNO Scintillator
  * Recent Opticks Enhancements : directed by Muon Production Experience
  * GPU Hit Merging : High Level Parallelization with CUDA Thrust
  * GPU Hit Merging : Avoids hiding Opticks performance

* Opticks + JUNOSW : Testing as Geometry and Models continue to Evolve

  * EMF Compensation Coils, Water Distributor, "Simtrace"


* Geant4 + Opticks + NVIDIA OptiX : Production Scaling ?

  * Geant4 + Opticks + NVIDIA OptiX : Monolith x4, x16
  * OpticksClients + OpticksService : Share GPUs 
  * Opticks "Optical Core" => Server, "Periphery" => Client
  * NP\_CURL.h : Array transport via HTTP POST 
  * Opticks Server Prototype : python + FastAPI + nanobind + CSGOptiX
  * NVIDIA Triton Inference Server (aka Dynamo-Triton)

* Summary + Links



POINTS FROM PRESENTATION PAGES TO COVER IN PROCEEDINGS
========================================================

Introduction
-------------


High Energy Adaption
---------------------

* p13: screenshot
* p16: results table presentation + benefits discussion
* p14: benefits of "lite" and merging

  * "lite" benefits + benefit twice from GPU hit merging
  * monster tests
  * async - future object (passing mention - if at all - as async so incomplete)

Testing as Geometry and Models continue to Evolve : NOT IN ABSTRACT - SKIP ?
---------------------------------------------------------------------------------

* p17: viz dev
* p18,19,20: small screenshots ?

Production Scaling ? NEED TO COVER BRIEFLY
--------------------------------------------

* p27: client/server rationale
* p28: NP\_CURL.h - just a mention
* p29: server prototype
* p30: Triton - dont mention vapourware 
* p31: GPU-less Opticks


IDEAS

* Intro OR Summary 
* find everyone referencing my papers - and reference them back with brief note on what progress they achieved  

}


%
\begin{document}
\title{Opticks: GPU Accelerated Optical Photon Simulation for JUNO and Other Experiments}
\author{\firstname{Simon} C. \lastname{Blyth}\inst{1}\fnsep\thanks{Corresponding author \email{simon.c.blyth@gmail.com}.}}
\institute{Institute of High Energy Physics, CAS, Beijing, China.}
\abstract{\input{opticks-blyth-chep2026-v1-abstract.tex}}
\maketitle
%
%
%
%
\section{Introduction}%
\label{intro}
%
Opticks[1-7] applies NVIDIA\textregistered\ OptiX\texttrademark~\cite{optix} GPU ray tracing
to optical photon simulation and integrates this with the Geant4 [8] simulation of all other particles.
Offloading optical photon simulation to the GPU eliminates processing time and memory bottlenecks especially for
large high-yield detectors such as the 20 kton Jiangmen Underground Neutrino Observatory (JUNO)~\cite{juno}. 
Although developed for JUNO, Opticks has been used by other experiments including LZ~\cite{lz}, NEXT~\cite{next}, DUNE~\cite{dune}
and the Electron Ion Collider~\cite{eic}. Prior CHEP proceedings [1–7] cover the development history of Opticks.
These proceedings focus on enhancements motivatated by JUNO production experience with the simulation
of high energy events, such as cosmic ray muons which can yield several billions
of optical photons per event.
%
%
\newpage
\subsection{Optical photon simulation problem}
Simulating optical photons resulting from cosmic ray muons crossing the JUNO liquid scintillator consumes
more than 99.5\% of total CPU time in standard Geant4~\cite{g4} based simulations. This time is dominated by 
ray tracing, calculating intersection positions of rays representing photon steps with the detector geometry,

Optical photon simulation is extremely well suited to GPU parallelization, due to
the simplicity of optical physics requiring minimal register memory per threads and massive
parallelism through billions of independent photons.  
NVIDIA OptiX~\cite{optix} provides a single-ray programming model that 
leverages state-of-the-art acceleration structures and dedicated RT Cores on NVIDIA RTX~\cite{rtx} GPUs
which has enabled Opticks to implement the full photon simulation lifecycle from generation through to
absorption or detection entirely in one kernel. 
%
%
%
%
\subsection{Hybrid CPU+GPU Simulation Workflow}
%
\label{secworkflow}%
%
\begin{figure}[t]
\centering
%                                    left lower right upper
\includegraphics[width=\textwidth,trim={0 4cm 0 4cm},clip]{env/Documents/Geant4OpticksWorkflow7/Geant4OpticksWorkflow7_005.png}
\caption{This figure is reproduced from~\cite{chep2023}, which provides further details. Illustration of the
hybrid Geant4 + Opticks workflow : {\tt G4CXOpticks} translates Geant4 geometry to GPU appropriate form. 
{\tt U4} collects "gensteps" enabling GPU generation of scintillation and Cerenkov photons. 
}
\label{figworkflow} 
\vspace{-5mm}
\end{figure}
%
%
Figure~\ref{figworkflow} shows the geometry and event worflow of Opticks.
At initialization the Geant4 volume tree and solids are translated into an NVIDIA OptiX geometry
and material and surface properties are collected into GPU textures, as detailed in prior proceedings. 
%
The Geant4 generation of optical photons secondary tracks via loops within the scintillation and Cerenkov processes
are replaced with the collection of "genstep" parameters, including the number of photons to generate, into a genstep array.
The gensteps are uploaded to the GPU, enabling an equivalent optical photon generation to be implemented
within an NVIDIA OptiX ray generation program.
%
Opticks CUDA launches conceptually generate and propagate all optical photons in parallel, with one CUDA thread per optical photon.
In practice the number of photons that can be simulated in parallel is constrained by the VRAM needed
to persist the configured photon device arrays. 
At initialization the total GPU VRAM is queried and a heuristic calculation of the maximum number of photon slots
that can be simulated in a single launch is estimated ({\tt SEventConfig::MaxSlot}).
Subsequently for each event a vector of slices of the genstep array {\tt SGenstep::GetGenstepSlices}
collected from Geant4 is formed such the photons from every genstep array slice stay within VRAM.
%








%
\newpage
\subsection{Optical Photon Simulation Implemented within OptiX Ray Generation Program}
\label{sec:kernel}
\begin{samepage}
%
Listing~\ref{lst:propagatekernel} summarizes the NVIDIA OptiX ray generation program
that implements the full simulation of a single optical photon from generation through
to absorption or detection. Scintillator reemission is treated as "rebirth" 
of the photon within the same single thread. Each optical photon thread has a unique {\tt photonIdx}
allowing all threads to simultaneously write into photon arrays.
%The seed array gives access for each photon to the associated parent genstep parameters.
%
The trace call at each while loop iteration
finds the intersection of a ray with the current photon position and momentum direction
with the detector geometry. The prd struct includes the distance to the intersection 
and the surface normal at intersection. The {\tt qsim::propagate} call implements
scattering, reemission, boundary physics, surface effects, absorption and detection.  
Optionally the final photon can be persisted into device global memory using
arrays of 64(16) byte {\tt sphoton(sphotonlite)} struct.
Also a "record" array can persist the photon history at up to maxRecord steps enabling
detailed debugging and animated photon history visualizations.   
%
\begin{lstlisting}[caption={Optical photon simulation implemented within NVIDIA OptiX ray generation program.}, label={lst:propagatekernel}]
sevent* evt = params.evt ;
unsigned genstepIdx = evt->seed[photonIdx] ;
const quad6& gs = evt->genstep[genstepIdx] ;

qsim* sim = params.sim ;

RNG rng ; // random number generation state 
sim->rng->init( rng, evt->index, photonIdx );

sphoton p ;
sim->generatePhoton(p, rng, gs, photonIdx, genstepIdx );

quad2 prd ; // per-ray-data

int command = START ;
int bounce = 0 ; 
while( bounce < params.maxBounce && p.time < params.maxTime )
{   
    trace(params.geom, p.pos, p.mom, params.tmin, params.tmax, prd );
    command = sim->propagate(bounce, rng, p);

    if(evt->record && bounce < evt->maxRecord) 
        evt->record[evt->maxRecord*photonIdx+bounce] = p ;

    bounce++;
    if(command == BREAK) break ; // photon absorbed or detected
}   

if( evt->photon ) evt->photon[photonIdx] = p ;   // final photon
if( evt->photonlite )
{   
    sphotonlite l ;
    sphotonlite::summarizePhoton(l, p) ; 
    evt->photonlite[photonIdx] = l ;
}   
\end{lstlisting}
\end{samepage}
%
\subsection{Opticks event arrays and simulation workflow}
%
\begin{center}
\begin{table}
\begin{tabular}{ |m{20mm}|m{20mm}|m{30mm}|m{40mm}| } 
 \hline
Event Arrays                 &  Struct/Type      &  NumPy shape               &  Note      \\ 
 \hline
   genstep                   & {\tt quad6}       &  (num\_genstep,6,4)        & Generation parameters \\
   seed                      & {\tt int}         &  (num\_photon)             & Associate photon to genstep \\
 \hline
   photon                    & {\tt sphoton}     &  (num\_photon,4,4)         & 64 byte details \\
   photonlite                & {\tt sphotonlite} &  (num\_photon,4)           & 16 byte summary  \\
 \hline
   record                    & {\tt sphoton}     &  (num\_photon,32,4,4)      & Photon step history \\  
   seq                       & {\tt sseq}        &  (num\_photon,4,2)         & Summary flag sequence \\  
 \hline
   hit                       & {\tt sphoton}     &  (num\_hit,4,4)            & Selection of photon \\
   hitmerged                 & {\tt sphoton}     &  (num\_hitmerged,4,4)      & Reduction of hit \\
   hitlite                   & {\tt sphotonlite} &  (num\_photon,4)           & Selection of photonlite \\
   hitlitemerged             & {\tt sphotonlite} &  (num\_hitmerged,4)        & Reduction of hitlite   \\
 \hline
\end{tabular}
\caption{\label{tabeventarrays}Principal event arrays used on device by the optical photon simulation. NumPy shapes use 32 bit elements.}
\end{table}
\vspace{-8mm}
\end{center}%
%
The main event arrays used by the Opticks GPU optical photon simulation are shown in Table~\ref{tabeventarrays}.
The genstep array data collected from Geant4 is uploaded to the device prior to launching the simulation kernel 
and the seed array is filled with genstep indices for each photon expanded from the photon counts for each genstep
using CUDA Thrust parallel primitives ({\tt sysrap/iexpand.h}).
The seed array enables the parent genstep for each photon slot to be accessed and the
photon to be generated. The simulation kernel described in section~\ref{sec:kernel} populates the photon(photonlite) arrays with the final photon state
at absorption or detection. In addition the simulation kernel can fill the "record" array with the full photon step history and
the "seq" array flag sequence with 4 bit flags from each step of the photon.
  
Following the simulation kernel typically one of the hit, hitmerged, hitlite or hitlitemerged arrays is populated on device
directly from the corresponding photon(photonlite) array. The implementation of hit selection and merging is described
in section~\ref{sec:hitmerge}. All arrays can be downloaded and persisted into NumPy format files for loading into python.
%
\section{Simulation of High Energy Events}
%
In large scale liquid scintillator experiments such as JUNO, high-energy events such
as cosmic ray muons can generate hundreds of millions to billions of optical photons.
Events with multiple muons with long paths crossing the liquid scintillator yield the most optical photons.  
With such events recording hits for every photon collected at photomultiplier tubes could
yield gigabytes of hit data of limited utility. Many studies do not require full details on 
optical photons collected at PMTs, thus the simulation provides an option to collect and save
only summary hit information. Using summary hits and merging hits arriving onto the same PMT within 
configurable time buckets into an aggregate record with the number of hits and the earliest hit time
allow greatly reduces data sizes.

A CPU hit merging algorithm based on collections of non-merged PMT hit pointers
for each PMT is invoked from the G4VSensitiveDetector::ProcessHits implementation used
within the standard JUNO simulation. For CPU simulation this approach is adequate as the
processing time for the merging is a small fraction of the overall processing time. 

Initially the Opticks integrated JUNO simulation reused this CPU hit merging algorithm, 
performing the merge of all hit photons downloaded from GPU buffers. The processing 
time for CPU merging following GPU simulation was however found to dominate the 
processing time for the optical photon simulation.  This situation motivated development
of GPU hit merging and a "lite" version of photon struct to allow Opticks to
profit from lower GPU resource usage when summary hits are configured. 
%
\subsection{CUDA Thrust Implementations of GPU hit selection and merging} 
\label{sec:hitmerge}
%
When GPU hit merging is not configured the hit(hitlite) arrays are
selected from the photon(photonlite) arrays with a device to device copy based on photon flag mask indicating detection 
using CUDA Thrust stream compaction via {\tt thrust::copy\_if}.
%
When hit merging is enabled the hitmerged(hitlitemerged) arrays are selected and merged from the
photon(photonlite) arrays with a single CUDA Thrust based method {\tt SPM::merge\_partial\_select}
that avoids intermediate arrays and additional processing steps. 
%
As hit selection and merging are applied to both the 64(16) byte {\tt sphoton(sphotonlite)} struct
arrays templated implementations are used with nested functors within the templated types
that specialize the selection, sorting and merging as described in table~\ref{tabfunctor}.
%
%
\begin{center}
\begin{table}
\begin{tabular}{ |m{20mm}|m{100mm}| }
\hline
Nested Functor          &  Action       \\
\hline\hline
{\tt key\_functor}      &  Returns identity integer for a hit within a time bucket  \\  
\hline
{\tt reduce\_op}        &  Merge two instances, adding hit counts and choosing earlier time  \\ 
\hline
{\tt select\_pred}      &  Selection predicate comparing photon flag mask with selection mask \\    
\hline
\end{tabular}
\caption{\label{tabfunctor}Nested functors of {\tt sphoton} and {\tt sphotonlite} that are used with the CUDA Thrust primitives to specialize selection, sorting and merging as
appropriate to the different structs.} 
\end{table}
\vspace{-8mm}
\end{center}%
%
%
Opticks optical photon simulation makes extensive use of the CUDA Thrust library of parallel
primitives which provides a high level way to benefit from GPU parallel performance.
The principal CUDA Thrust C++ primitives used to implement GPU hit selection and 
merging are summarized in Table~\ref{thrustmethods}.

For events where VRAM is sufficient to simulate all photons in a single launch a single call to the select and merge method
yields the configured hit/hitmerged/hitlite/hitlitemerged array. Events requiring multiple launches are merged
after each launch with the resulting arrays concatenated on the CPU and uploaded to the device prior to
performing a final merge using the same method with a special cased flagmask argument that indicates 
that no selection is done. 
%
\begin{center}
\begin{table}
\begin{tabular}{ |m{80mm}|m{30mm}| } 
 \hline
 Action                                & CUDA Thrust method                \\ 
\hline\hline 
  Select hits from photons             & {\tt copy\_if}         \\ 
\hline 
  Form key(pmtid,timebucket) from hit  & {\tt transform}        \\ 
\hline 
  Sort hits by hash(pmtid,timebucket)  & {\tt sort\_by\_key}    \\ 
\hline 
  Merge hit: earlier time, sum counts  & {\tt reduce\_by\_key} \\ 
 \hline
\end{tabular}
\caption{\label{tabthrustmethods}Principal CUDA Thrust primitives used to implement GPU hit merging. }
\end{table}
\vspace{-8mm}
\end{center}%


Table~\ref{tab:opticks_junosw_performance} shows performance in four modes named after the type of hit arrays that are downloaded from GPU.
The merged modes apply hit merging on the GPU avoiding post processing on the CPU. As the GPU ray trace accelerated simulation is exactly
the same in the four cases the kernel simulation time is approximately the same. 

The 6.3 seconds time difference between hitmerged and hitlitemerged modes is dominated by hit localization that is not required for lite photons, as
in that case this localization is done during the simulation. The localization applies 4x4 transforms to every global coordinate hit to give positions within the 
frame of the PMT is currently performed on the CPU. Implementing GPU hit localization is expected to remove this time difference.

Performing the hit merging on the GPU is doubly beneficial as it eliminates the need for slow CPU merging and it
reduces the data downloaded by a factor of four due to the reduced number of merged hits.
In these measurements the benefits from using "lite" photon are not visible as 150M photons were simulated in a single launch.
With 32GB of VRAM the maximum number of photons that can be simulated per launch is approximately 250M with full 64 byte photons
and 500M with 16 byte photons. This max photon increase is only a factor of two, not a factor of four that would be suggested by struct sizes, 
due to transient VRAM usage needed by CUDA Thrust for the hit merging. Simulation of larger events requiring two or more launches for full photons
but only one for lite photons are needed to quantify the advantages from the lite photons and also the advantages of using more VRAM.
The benefits from using lite photons might be considered equivalent to the benefit of doubling available VRAM.
%
\begin{table}[htbp]
  \centering
  \small
  \begin{tabular}{llccccc}
    \toprule
    \textbf{Opticks+JSW}    & \textbf{merge} & \textbf{Simulate}     & \textbf{(GPU merge)+}   & \textbf{(CPU merge)+}    & \textbf{Total [s]} & \textbf{Speedup} \\
    \textbf{hit\_mode}      &                & \textbf{Kernel [s]}   & \textbf{Download [s]}   & \textbf{Collection [s]}  &                    & \textbf{vs std}  \\
    \midrule
    hit           & \multirow{2}{*}{CPU}     & 23.0                  & 1.95                    & 190.4                    & 215.6              & x32             \\
    hitlite       &                          & 23.1                  & 0.48                    & 146.5                    & 170.2              & x31             \\
    \midrule                                                                                                                                   
    hitmerged     & \multirow{2}{*}{GPU}     & 23.0                  & \btext{0.54}            & 6.7                      & 30.4               & \rtext{x233}    \\
    hitlitemerged &                          & 23.1                  & \btext{0.18}            & 0.4                      & 23.8               & \rtext{x221}    \\
    \bottomrule
  \end{tabular}
  \caption{\label{tab:opticks_junosw_performance}Comparison of processing times in various hit modes and speedups relative to standard single threaded Geant4 with the JUNO simulation (JSW),
for one double muon event with 150M photons, 28M hit and 6.4M mergedHit using and 1 ns merge bucket. The "lite" modes write {\tt sphotonlite}. 
Timings obtained using Dell Precision 7960 workstation with NVIDIA RTX 5000 Ada generation GPU (third generation RT cores) with 32 GB of VRAM.}
\end{table}
%
%
%
\section{Scaling Opticks for Monte Carlo Production}

Opticks eliminates the optical photon bottleneck in Geant4 simulations;
however, sustaining high GPU resource utilization across production Monte Carlo campaigns 
with diverse workloads and hardware capabilities remains challenging.
Optical simulation performance has roughly doubled with each NVIDIA RTX GPU generation 
due to ray-tracing hardware advances, making GPU starvation an increasingly critical issue.
In a monolithic Opticks application, the GPU remains idle during Geant4 initialization 
and non-optical simulation phases. 
While overlapping CPU generation with GPU propagation can yield moderate gains, 
adopting a client-server architecture offers the architectural flexibility to fully saturate GPU resources,
eliminate redundant geometry initialization costs, and mitigate GPU hardware scarcity. 
Crucially, tuning the ratio of CPU clients to GPU servers in production campaigns affords fine-grained 
control over GPU utilization, decoupled from specific workloads or hardware capbilities.

The modular package structure of Opticks enables a clean division into client and server components 
compiled from a single codebase. Core optical simulation packages requiring CUDA and OptiX dependencies 
are isolated on the remote server, which services requests from lightweight clients. 
The current prototype server uses Python, FastAPI, and nanobind to expose the C++ \texttt{CSGOptiX} simulation engine, 
with the client relying on \texttt{libcurl} to provide HTTP transport.

Future work will explore replacing the prototype with the NVIDIA Triton Inference Server using 
gRPC transport and Protocol Buffers. This may provide more robust multi-client performance
and higher throughput when returning large hit arrays.



\section{Summary}
%
%Opticks enables Geant4-based optical photon simulations to benefit from 
%state-of-the-art NVIDIA GPU ray tracing, made accessible via the NVIDIA OptiX 7+ API,
%allowing memory and time processing bottlenecks to be eliminated. 
%Recent feature additions such as out-of-core optical photon simulation and the 
%adoption of the curand Philox random number generator remove limits on the number of 
%photons that can be simulated and make Opticks easier to use. 
%Comparisons between different GPUs shows optical simulation performance 
%to closely follow ray tracing performance with a factor of four improvement between 
%GPUs from the first and third RTX generations. 


Shader Execution Reordering (SER) introduced with third RTX generation GPUs has
been reported~\cite{ser} with the Simphony~\cite{simphony} Opticks fork to achieve more
than a factor of 10 reduction in propagation kernel times by mitigating warp divergence.




%
%
%\newpage
\section*{Acknowledgements}
%
The JUNO collaboration is acknowledged for the use of detector 
geometries and simulation software. Dr. Tao Lin is acknowledged 
for his JUNO software assistance over many years. NVIDIA is acknowledged for extensive 
support for the OptiX 7+ API transition.
%
This work is supported by National Natural Science Foundation of China (NSFC)
under grant No. 12275293.
%
\begin{thebibliography}{}
%
%1 
\bibitem{opticks}
Opticks Repository, {\tt https://bitbucket.org/simoncblyth/opticks/}\\
Opticks References, {\tt https://simoncblyth.bitbucket.io}\\
Opticks Group, {\tt https://groups.io/g/opticks}
%
%2
\bibitem{chep2024}
S. C. Blyth, EPJ Web Conf. {\bf 337}, 11093 (2025) \\
{\tt https://doi.org/10.1051/epjconf/202533701093}
%
\bibitem{chep2023}
S. Blyth, EPJ Web Conf. {\bf 295}, 11014 (2024) \\
{\tt https://doi.org/10.1051/epjconf/202429511014}
%3
\bibitem{chep2021}
S. Blyth, EPJ Web Conf. {\bf 251}, 03009 (2021) \\
{\tt https://doi.org/10.1051/epjconf/202125103009}
%4
\bibitem{chep2019}
S. Blyth, EPJ Web Conf. {\bf 245}, 11003 (2020) \\
{\tt https://doi.org/10.1051/epjconf/202024511003}
%5
\bibitem{chep2018}
S. Blyth, EPJ Web Conf. {\bf 214}, 02027 (2019) \\
{\tt https://doi.org/10.1051/epjconf/201921402027}
%6 
\bibitem{chep2016}
Blyth Simon C 2017 J. Phys.: Conf. Ser. {\bf 898} 042001 \\
{\tt https://doi.org/10.1088/1742-6596/898/4/042001}
%
%7
\bibitem{g4}
S. Agostinelli et al., Nucl. Instrum. Methods. Phys. Res. A {\bf 506}, 250 (2003)\\
\hspace*{0.0em}J. Allison et al., IEEE Trans Nucl Sci, {\bf 53}, 270 (2006)\\
\hspace*{0.0em}J. Allison et al., Nucl. Instrum. Methods. Phys. Res. A {\bf 835}, 186 (2016)
%
%   Recent Developments in Geant4, J. Allison et al., Nucl. Instrum. Meth. A 835 (2016) 186-225
%   Geant4 Developments and Applications, J. Allison et al., IEEE Trans. Nucl. Sci. 53 (2006) 270-278
%   Geant4 - A Simulation Toolkit, S. Agostinelli et al., Nucl. Instrum. Meth. A 506 (2003) 250-303
%
%8
\bibitem{optix}
S. Parker et al., ACM Trans. Graph.: Conf. Series {\bf 29}, 66 (2010)\\
\hspace*{0.0em}OptiX introduction, {\tt https://developer.nvidia.com/optix}\\
\hspace*{0.0em}OptiX API, {\tt https://raytracing-docs.nvidia.com/optix8/index.html}
%
%9
\bibitem{rtx}
NVIDIA RTX\texttrademark\, Platform, {\tt https://developer.nvidia.com/rtx}
%10
\bibitem{juno}
Neutrino physics with JUNO \\
F. An et al., J. Phys. G. {\bf 43}, 030401 (2016) 
%
%
\bibitem{next}
NEXT Collaboration. Performance of an optical TPC Geant4 simulation with opticks GPU-accelerated photon propagation.\\
Eur. Phys. J. C 85, 910 (2025). https://doi.org/10.1140/epjc/s10052-025-14612-0
%
\bibitem{dune}
Accelerating Optical Photon Simulation in DUNE with Opticks\\
Ilker Parmaksiz, Aaron Higuera, Laura Paulucci, Viktor Pec, Estanislao Forino\\
arXiv:2608.27306 [hep-ex]\\
DOI: https://doi.org/10.48550/arXiv.2608.27306
%
\bibitem{lz}
GPU simulation with Opticks: The future of optical simulations for LZ \\
Oisín Creaner, Simon Blyth, Sam Eriksen, Lisa Gerhardt, Maria Elena Monzani and Quentin Riffard\\
EPJ Web Conf., 251 (2021) 03037\\
DOI: https://doi.org/10.1051/epjconf/202125103037
%
\bibitem{eic}
GPU acceleration of optical photon propagation in low photon yield applications: Opticks for the Electron Ion Collider.\\
Galgoczi, G., Kauder, K., Potekhin, M. et al.\\
Eur. Phys. J. C 86, 934 (2026).\\
https://doi.org/10.1140/epjc/s10052-026-16180-3
%
\bibitem{ser}
Mitigating warp divergence in GPU optical photon Monte Carlo: an order of magnitude speedup with Shader Execution Reordering.\\
Galgoczi, G. (2026), arXiv:2608.21396 [physics.ins-det]. https://arxiv.org/abs/2608.21396
%
\bibitem{simphony}
G. Galgoczi et al., GPU optical photon Monte Carlo for noble liquid detectors:
validation against Geant4 in a large liquid argon TPC benchmark, arXiv:2606.05385
(2026).
%\bibitem{amdahl}
%Validity of the single processor approach to achieving large scale computing capabilities.\\
%Amdahl, G. M. (1967).\\ 
%Proceedings of the April 18-20, 1967, Spring Joint Computer Conference, 483–485. https://doi.org/10.1145/1465482.1465560\\
%
%11
\bibitem{throughput}
Understanding Throughput Oriented Architectures \\
M. Garland, D.B. Kirk, Commun. ACM {\bf 53}(11), 58 (2010) 
%12
\bibitem{numpy}
The NumPy array: a structure for efficient numerical computation \\
S. Van der Walt, S. Colbert, G. Varoquaux, Comput. Sci. Eng. {\bf 13}, 22 (2011)
%
%13
\bibitem{curand}
cuRAND, {\tt http://docs.nvidia.com/cuda/curand/index.html}\\
\hspace*{0.0em}Generator tests, {\tt https://docs.nvidia.com/cuda/curand/testing.html}\\
\hspace*{0.0em}{\tt https://docs.nvidia.com/cuda/curand/device-api-overview.html\#performance-notes}
%
%
\end{thebibliography}
%
\end{document}
